\subsection{profinite 纤维扩张、局部化与缠绕度有理化}\label{subsec:cdim-profinite-fiber-localization-rational-winding}
本小节在定义 \ref{def:circle-dimension} 与定义 \ref{def:half-circle-dimension} 的基础上，将圆维口径从有限生成离散交换群扩展到一般离散交换群，并把 prime-register 的 profinite 纤维结构与 solenoid 自同态的度值域统一为同一条可审计链。
核心目标是分离两类信息：连续相位圆因子的维数（主分量）与素数纤维坐标（签名分量）。

\begin{definition}[扩展圆维]\label{def:cdim-star}
设 \(A\) 为离散交换群，\(\widehat A:=\operatorname{Hom}(A,\TT)\) 为其 Pontryagin 对偶紧交换群，\((\widehat A)^0\) 为单位元连通分量。
定义
\[
\cdim^\ast(A):=\dim_{\mathrm{top}}\!\bigl((\widehat A)^0\bigr),
\]
其中 \(\dim_{\mathrm{top}}\) 为紧 Hausdorff 空间的覆盖维数。
\end{definition}

\begin{proposition}[与有限生成口径的一致性]\label{prop:cdim-star-compatible-fg}
若 \(A\) 为有限生成离散交换群，则
\[
\cdim^\ast(A)=\cdim(A).
\]
\end{proposition}
\begin{proof}
由有限生成交换群结构定理，\((\widehat A)^0\) 为有限维环面，覆盖维数与其 Lie 维数一致。
故定义 \ref{def:cdim-star} 退化为定义 \ref{def:circle-dimension}。证毕。
\leanverified{paper\_cdim\_star\_compatible\_fg\_seeds}
\leanverified{paper\_cdim\_star\_compatible\_fg\_package}
\leanverified{paper\_cdim\_star\_compatible\_fg}
\end{proof}

\begin{theorem}[torsion/profinite 纤维扩张不增扩展圆维]\label{thm:cdim-star-torsion-fiber-invariance}
\leanverified{paper\_cdim\_star\_torsion\_fiber\_invariance}
设
\[
0\longrightarrow F\longrightarrow G\longrightarrow H\longrightarrow 0
\]
为离散交换群短正合列，且 \(F\) 为 torsion 群。
则
\[
(\widehat G)^0\cong (\widehat H)^0,
\qquad
\cdim^\ast(G)=\cdim^\ast(H).
\]
\end{theorem}
\begin{proof}
对偶化得短正合列
\[
0\longrightarrow \widehat H\longrightarrow \widehat G\longrightarrow \widehat F\longrightarrow 0.
\]
由于 \(F\) 为 torsion，\(\widehat F\) 为紧且全不连通群。
于是 \((\widehat G)^0\) 在商群 \(\widehat G/\widehat H\cong \widehat F\) 中的像是连通子群，只能是单位元，故
\[
(\widehat G)^0\subseteq \widehat H.
\]
左侧为连通子群，故必包含于 \((\widehat H)^0\)，即 \((\widehat G)^0\subseteq (\widehat H)^0\)。
反向包含由 \(\widehat H\subseteq \widehat G\) 立即得到：\((\widehat H)^0\subseteq (\widehat G)^0\)。
两者相等，维数相等。证毕。
\end{proof}

\begin{theorem}[连通相位与 profinite 纤维的不可分解性]\label{thm:cdim-connected-phase-profinite-fiber-indecomposable}
\leanverified{paper\_cdim\_connected\_phase\_profinite\_fiber\_indecomposable}
设 \(r\ge 1\) 并给定紧交换群短正合列
\[
0\longrightarrow P\longrightarrow G\xrightarrow{\ \pi\ }\TT^r\longrightarrow 0,
\]
其中 \(G\) 连通而 \(P\) 为非平凡 profinite 紧交换群（因而全不连通）。
则该扩张在拓扑群范畴中不分裂；等价地，不存在拓扑群同构 \(G\cong \TT^r\times P\)。
更强地，任意连续群同态 \(f:G\to P\) 必为平凡同态。
\end{theorem}
\begin{proof}
紧连通群的连续像仍连通。
由于 \(P\) 为全不连通群，其连通子群只能为单位元，故任意连续 \(f:G\to P\) 的像为 \(\{0\}\)，从而 \(f=0\)。
\par
若扩张分裂，则存在连续截面 \(s:\TT^r\to G\) 使 \(\pi\circ s=\mathrm{id}\)，从而 \(G\cong \TT^r\times P\)。
但当 \(P\neq 0\) 时，\(\TT^r\times P\) 非连通，与 \(G\) 连通矛盾。证毕。
\end{proof}

\begin{theorem}[局部化群 \texorpdfstring{$\ZZ[1/S]$}{Z[1/S]} 的对偶扩张与一维性]\label{thm:cdim-star-z1s-dual-extension}
设 \(S\) 为有限素数集，记
\[
A_S:=\ZZ[1/S].
\]
则存在紧交换群短正合列
\[
0\longrightarrow \prod_{p\in S}\ZZ_p
\longrightarrow \widehat{A_S}
\longrightarrow \TT
\longrightarrow 0,
\]
并且
\[
\cdim^\ast(A_S)=1.
\]
\leanverified{paper\_cdim\_star\_z1s\_dual\_extension}
\end{theorem}
\begin{proof}
由短正合列
\[
0\to \ZZ\to A_S\to A_S/\ZZ\to 0
\]
及 \(p\)-primary 分解，有
\[
A_S/\ZZ\cong \bigoplus_{p\in S}\ZZ(p^\infty),
\]
其中 \(\ZZ(p^\infty)\) 为 Prüfer \(p\)-群。
对偶化并使用
\(
\widehat{\ZZ(p^\infty)}\cong \ZZ_p
\)
与 \(\widehat{\ZZ}\cong\TT\)，得所示短正合列。

另一方面，\(A_S\) 为 rank-\(1\) torsion-free 群，可写成 \(\ZZ\xrightarrow{\times n_1}\ZZ\xrightarrow{\times n_2}\cdots\) 的直极限（其中每个 \(n_j\) 的素因子属于 \(S\)），故
\[
\widehat{A_S}\cong \varprojlim\!\bigl(\TT\xleftarrow{z\mapsto z^{n_1}}\TT\xleftarrow{z\mapsto z^{n_2}}\cdots\bigr).
\]
该逆极限由一维紧空间构成，故覆盖维数不超过 \(1\)；同时其到任一阶段圆周的投影为满射，故维数至少 \(1\)。
因此 \(\dim_{\mathrm{top}}(\widehat{A_S})=1\)，即 \(\cdim^\ast(A_S)=1\)。证毕。
\end{proof}

\begin{corollary}[procyclic profinite 纤维的一圆维 solenoid 分类]\label{cor:cdim-procyclic-fiber-solenoid-classification}
\leanverified{paper\_cdim\_procyclic\_fiber\_solenoid\_classification}
设 \(K\) 为 procyclic 的 profinite 紧交换群，且存在短正合列
\[
0\longrightarrow K\longrightarrow G\longrightarrow \TT\longrightarrow 0,
\]
其中 \(G\) 为紧连通交换群。
则存在有限素数集 \(S\) 使得
\[
G\ \cong\ \widehat{\ZZ[1/S]},
\qquad
K\ \cong\ \prod_{p\in S}\ZZ_p.
\]
若进一步令 \(N:=\prod_{p\in S}p\)，则亦可写为
\(
G\cong \widehat{\ZZ[1/N]}
\)
且
\(
K\cong \ZZ_N\cong \prod_{p\mid N}\ZZ_p
\)；
其中素数集合 \(S=\{p:\mathrm{Sylow}_p(K)\neq 0\}\) 由 \(K\) 的 Sylow 结构唯一确定。
\end{corollary}
\begin{proof}
对偶化给出短正合列
\[
0\longrightarrow \ZZ\longrightarrow A\longrightarrow K^\vee\longrightarrow 0,
\qquad A:=\widehat G,
\]
其中 \(A\) 为秩 \(1\) 的无挠离散群而 \(K^\vee\) 为挠群。
由于 \(K\) procyclic，\(K^\vee\) 为有限个 Prüfer 群 \(\ZZ(p^\infty)\) 的直和；记其素数支持为 \(S\)。
于是 \(A/\ZZ\cong \bigoplus_{p\in S}\ZZ(p^\infty)\)。
而对秩一无挠子群 \(A\)（且含 \(\ZZ\)），该等价于 \(A\cong \ZZ[1/S]\)。
再对偶化即得 \(G\cong \widehat{\ZZ[1/S]}\) 与
\(
K\cong \widehat{\bigoplus_{p\in S}\ZZ(p^\infty)}\cong \prod_{p\in S}\ZZ_p
\)。
最后的唯一性由 procyclic profinite 群的 Sylow 分解与 \(\ZZ_p\) 的正交性（定理 \ref{thm:cdim-star-prime-coordinate-orthogonality}）给出。证毕。
\end{proof}

\begin{corollary}[有理群的相位圆与全素数纤维分解]\label{cor:cdim-star-q-dual-extension}
\leanverified{paper\_cdim\_star\_q\_dual\_extension}
对离散群 \((\QQ,+)\)，有短正合列
\[
0\longrightarrow \widehat{\ZZ}\cong\prod_{p}\ZZ_p
\longrightarrow \widehat{\QQ}
\longrightarrow \TT
\longrightarrow 0,
\]
并且
\[
\cdim^\ast(\QQ)=1.
\]
\end{corollary}
\begin{proof}
由
\(
\QQ/\ZZ\cong\bigoplus_p \ZZ(p^\infty)
\)
对偶化即得核 \(\prod_p\ZZ_p\) 与商 \(\TT\) 的短正合列。
同理于定理 \ref{thm:cdim-star-z1s-dual-extension} 的逆极限口径，\(\widehat\QQ\) 为一维通用 solenoid，故 \(\cdim^\ast(\QQ)=1\)。证毕。
\end{proof}

\begin{remark}[完备相位对象并非裸圆]\label{rem:cdim-complete-phase-object-circle-plus-fiber}
定理 \ref{thm:cdim-star-z1s-dual-extension} 与推论 \ref{cor:cdim-star-q-dual-extension} 表明，秩一局部化与有理完成的可见连通商始终是 \(\TT\)，但完整对象同时携带全不连通的 profinite 素数纤维。因而圆在本文中承担的是可见相位商的角色，而非完备对象的全部内容。
\par
换言之，\(\widehat{\ZZ[1/S]}\) 与 \(\widehat\QQ\) 的正确理解均不是“单个裸圆”，而是“圆商与零维素数纤维的统一系统”。忽略后者，只能得到 Lie 可见层的投影，而得不到可逆、可复核且对完成化稳定的全对象。
\end{remark}

\begin{proposition}[prime-register 轴的实现层注入与有限深度塌缩]\label{prop:cdim-profinite-prime-register-embedding-collapse}
\leanverified{paper\_cdim\_profinite\_prime\_register\_embedding\_collapse}
记 profinite 整数环
\[
\widehat{\ZZ}:=\varprojlim_{N}\ZZ/N\ZZ\ \cong\ \prod_{p}\ZZ_p.
\]
则对任意 \(n\ge 2\)，存在显式单射群同态
\[
\iota_n:\ZZ^n\hookrightarrow \widehat{\ZZ},
\qquad
\iota_n(a_0,\dots,a_{n-1})
:=
\Bigl(\sum_{i=0}^{n-1}a_i\,p^{\,i}\Bigr)_{p\ \mathrm{prime}}.
\]
并且对任意 \(N\ge 2\)，合成到有限商环的映射
\[
\ZZ^n\xrightarrow{\ \iota_n\ }\widehat{\ZZ}\twoheadrightarrow \ZZ/N\ZZ
\]
不可能为单射；更强地，存在 \(0\neq v\in\ZZ^n\) 使 \(\iota_n(v)\equiv 0\pmod N\)。
\end{proposition}
\begin{proof}
\textbf{单射性.}
设 \(\iota_n(a_0,\dots,a_{n-1})=0\)。
则对每个素数 \(p\) 都有
\[
\sum_{i=0}^{n-1}a_i\,p^{\,i}=0\quad\text{在 }\ZZ_p\text{ 中}.
\]
左端为整数，若其在 \(\ZZ_p\) 中为零，则该整数本身为零。
因此多项式
\(
F(T):=\sum_{i=0}^{n-1}a_iT^i\in\ZZ[T]
\)
在无穷多个整数点 \(T=p\)（素数）处取零值，只能推出 \(F\equiv 0\)，从而 \(a_i=0\) 对一切 \(i\) 成立。
\par
\textbf{有限深度塌缩.}
自然投影 \(\widehat{\ZZ}\twoheadrightarrow \ZZ/N\ZZ\) 由逆极限定义给出。
其像至多含 \(N\) 个元素，而集合 \(\{0,1,\dots,N-1\}^n\subset\ZZ^n\) 含 \(N^n>N\) 个元素（因 \(n\ge 2\)）。
由鸽巢原理，存在 \(u\neq v\) 使二者在 \(\ZZ/N\ZZ\) 中同像，于是 \(w=u-v\neq 0\) 满足 \(\iota_n(w)\equiv 0\pmod N\)。证毕。
\end{proof}

\begin{theorem}[混合局部化的相位圆轴与素数纤维轴分解]\label{thm:cdim-star-mixed-localization-two-axis-decomposition}
\leanverified{paper\_cdim\_star\_mixed\_localization\_two\_axis\_decomposition}
令 \(r\ge 1\)。
对每个 \(i\in\{1,\dots,r\}\) 取有限素数集 \(S_i\subset\mathbb P\)，并记
\[
A_i:=\ZZ[1/S_i],
\qquad
A:=\bigoplus_{i=1}^{r}A_i.
\]
令 \(\widehat A\) 为 \(A\) 的 Pontryagin 对偶紧交换群。
则存在自然短正合列
\[
0\longrightarrow K\longrightarrow \widehat A\xrightarrow{\ \pi\ }\TT^{r}\longrightarrow 0,
\]
且核具有显式素数分解
\[
K\ \cong\ \prod_{p\in\mathbb P}\ZZ_p^{\,m_p},
\qquad
m_p:=\card{\{\,i:\ p\in S_i\,\}}.
\]
从而
\[
\cdim^\ast(A)=r,
\qquad
\pcdim(K)=\max_{p\in\mathbb P}m_p,
\]
其中 \(\pcdim\) 为 pro-圆维（定义 \ref{def:cdim-pro-circle-dimension}）。
\end{theorem}
\begin{proof}
由 Pontryagin 对偶把离散直和送到紧直积，得到 \(\widehat A\cong \prod_{i=1}^{r}\widehat{A_i}\)。
对每个 \(i\)，定理 \ref{thm:cdim-star-z1s-dual-extension} 给出短正合列
\[
0\to \prod_{p\in S_i}\ZZ_p\to \widehat{A_i}\to \TT\to 0.
\]
对 \(i=1,\dots,r\) 取有限直积并用正合性在有限直积下保持，得
\[
0\to \prod_{i=1}^{r}\prod_{p\in S_i}\ZZ_p\to \widehat A\to \TT^{r}\to 0.
\]
取 \(K:=\prod_{i=1}^{r}\prod_{p\in S_i}\ZZ_p\) 即得第一条短正合列。
将有限指标的直积按素数坐标重排，得到
\(
K\cong\prod_{p\in\mathbb P}\ZZ_p^{\,m_p}
\)。
\par
由于 \(\TT^{r}\) 连通且 \(K\) 全不连通，\(\pi\) 将 \((\widehat A)^0\) 满射到 \(\TT^{r}\)，并且 \((K)^0=\{0\}\) 蕴含 \((\widehat A)^0\cong \TT^{r}\)。
由定义 \ref{def:cdim-star} 得 \(\cdim^\ast(A)=r\)。
\par
最后，由命题 \ref{prop:cdim-pro-cdim-crt-reuse}(2)，\(\pcdim(K)=\max_p\pcdim(K_p)\)，其中 \(K_p\cong \ZZ_p^{m_p}\)。
而 \(\ZZ_p^{m_p}\) 的最小拓扑生成元数为 \(m_p\)，故 \(\pcdim(K)=\max_p m_p\)。证毕。
\end{proof}

\subsubsection{有限局部化账本的精确序列、乘法线性化障碍与商嵌不对称}\label{subsubsec:cdim-finite-localization-ledger-gap-rigidity}
在定理 \ref{thm:cdim-star-mixed-localization-two-axis-decomposition} 已把有限局部化直和的对偶结构压缩为 \(\TT^r\) 上的紧扩张之后，还可把这一对象族在离散侧、实现侧与拓扑侧进一步闭合为下列四条结论：有限局部化账本具有完全显式的离散--紧双短正合列；全部自然数乘法不能被任何有限秩局部化加性账本无损线性化；实现层最小编码与结构层有限秩圆维之间出现可任意放大的精确裂隙；而相应 solenoid 总可向有限维环面压缩，却在任何非平凡素数账本存在时都不能反向嵌入有限维环面。

\begin{theorem}[有限局部化直和的素数账本精确序列与有限秩圆维]\label{thm:cdim-finite-localization-directsum-prime-ledger-exact-sequence}
\leanverified{paper\_cdim\_finite\_localization\_directsum\_prime\_ledger\_exact\_sequence}
设
\[
\mathbf S=(S_1,\dots,S_r),
\qquad
S_i\subset\mathbb P\ \text{有限},
\]
并定义
\[
G_{\mathbf S}:=\bigoplus_{i=1}^{r}G_{S_i},
\qquad
G_{S_i}:=\ZZ[S_i^{-1}],
\qquad
m_p(\mathbf S):=\#\{\,i:\ p\in S_i\,\}.
\]
则存在自然短正合列
\[
0\longrightarrow \ZZ^r\longrightarrow G_{\mathbf S}
\longrightarrow \bigoplus_{p\in\mathbb P}C_{p^\infty}^{\,m_p(\mathbf S)}
\longrightarrow 0,
\]
其 Pontryagin 对偶为
\[
0\longrightarrow \prod_{p\in\mathbb P}\ZZ_p^{\,m_p(\mathbf S)}
\longrightarrow \widehat{G_{\mathbf S}}
\longrightarrow \TT^r
\longrightarrow 0.
\]
并且
\[
\cdim_{\mathrm{fr}}(G_{\mathbf S})=r.
\]
\end{theorem}
\begin{proof}
对每个 \(i\)，定理 \ref{thm:killo-localization-circle-dimension-prime-ledger-splitting} 给出短正合列
\[
0\longrightarrow \ZZ\longrightarrow G_{S_i}
\longrightarrow \bigoplus_{p\in S_i}C_{p^\infty}
\longrightarrow 0.
\]
对 \(i=1,\dots,r\) 取有限直和，得到
\[
0\longrightarrow \ZZ^r\longrightarrow G_{\mathbf S}
\longrightarrow \bigoplus_{i=1}^{r}\bigoplus_{p\in S_i}C_{p^\infty}
\longrightarrow 0.
\]
再按素数重排右端，即得
\[
\bigoplus_{i=1}^{r}\bigoplus_{p\in S_i}C_{p^\infty}
\cong
\bigoplus_{p\in\mathbb P}C_{p^\infty}^{\,m_p(\mathbf S)}.
\]
这就得到离散侧短正合列。

对其作 Pontryagin 对偶，并使用
\[
\widehat{\ZZ}\cong\TT,
\qquad
\widehat{C_{p^\infty}}\cong \ZZ_p,
\]
即可得到紧侧短正合列
\[
0\longrightarrow \prod_{p\in\mathbb P}\ZZ_p^{\,m_p(\mathbf S)}
\longrightarrow \widehat{G_{\mathbf S}}
\longrightarrow \TT^r
\longrightarrow 0.
\]

最后，每个 \(G_{S_i}\) 都是秩 \(1\) 的无挠阿贝尔群，故
\[
G_{\mathbf S}\otimes_{\ZZ}\QQ
\cong
\bigoplus_{i=1}^{r}\bigl(G_{S_i}\otimes_{\ZZ}\QQ\bigr)
\cong
\QQ^r.
\]
由定理 \ref{thm:conclusion-cdim-finite-rank-extension}，
\[
\cdim_{\mathrm{fr}}(G_{\mathbf S})
=
\operatorname{rank}_{\QQ}(G_{\mathbf S})
=
r.
\]
证毕。
\end{proof}

\begin{theorem}[任意有限局部化账本都不能同态化全部乘法]\label{thm:cdim-no-global-multiplicative-linearization-into-finite-localization-ledger}
\leanverified{paper\_cdim\_no\_global\_multiplicative\_linearization\_into\_finite\_localization\_ledger}
设 \(a,r\ge 0\)，\(\mathbf S=(S_1,\dots,S_r)\) 为有限素数集族，并记
\[
A_{a,\mathbf S}:=\ZZ^{a}\oplus G_{\mathbf S}.
\]
则不存在单射幺半群同态
\[
\Phi:\NN_{\times}\hookrightarrow A_{a,\mathbf S}
\]
满足
\[
\Phi(mn)=\Phi(m)+\Phi(n)\qquad(\forall\,m,n\in\NN_{\times}),
\]
其中 \(\NN_{\times}:=(\NN_{\ge 1},\times)\)。
\end{theorem}
\begin{proof}
反设存在这样的 \(\Phi\)。由命题 \ref{prop:killo-grothendieck-completion-preserves-injection}，对可消交换幺半群取 Grothendieck 完成后仍得到群单射
\[
G(\NN_{\times})\hookrightarrow A_{a,\mathbf S}.
\]
而
\[
G(\NN_{\times})
\cong
\QQ_{>0}^{\times}
\cong
\bigoplus_{p\in\mathbb P}\ZZ,
\]
为可数无穷秩自由阿贝尔群。

另一方面，由定理 \ref{thm:cdim-finite-localization-directsum-prime-ledger-exact-sequence}，
\[
\operatorname{rank}_{\QQ}(G_{\mathbf S})=r.
\]
故
\[
\operatorname{rank}_{\QQ}(A_{a,\mathbf S})=a+r<\infty.
\]
于是 \(A_{a,\mathbf S}\) 是有限秩无挠阿贝尔群。对任意阿贝尔群单射
\[
\bigoplus_{p\in\mathbb P}\ZZ\hookrightarrow A_{a,\mathbf S}
\]
张量 \(\QQ\) 后仍给出 \(\QQ\)-向量空间单射
\[
\bigoplus_{p\in\mathbb P}\QQ\hookrightarrow A_{a,\mathbf S}\otimes_{\ZZ}\QQ.
\]
左端维数可数无穷，右端维数仅为 \(a+r\)，矛盾。故所求 \(\Phi\) 不存在。证毕。
\end{proof}

\leanverified{paper\_cdim\_implementation\_structural\_gap\_for\_finite\_localization\_ledgers}
\begin{corollary}[实现层半圆维与结构层圆维的精确裂隙]\label{cor:cdim-implementation-structural-gap-for-finite-localization-ledgers}
对任意 \(a,r\ge 0\) 且 \(a+r>0\)，定义
\[
A_{a,\mathbf S}:=\ZZ^{a}\oplus G_{\mathbf S},
\]
并记
\[
\cdim^{\mathrm{impl}}(X)
:=
\inf\left\{\frac{m}{2}:\ \exists\,\text{集合单射 }X\hookrightarrow \NN^{m}\right\}.
\]
则
\[
\cdim^{\mathrm{impl}}(A_{a,\mathbf S})=\frac12,
\qquad
\cdim_{\mathrm{fr}}(A_{a,\mathbf S})=a+r,
\]
从而
\[
\cdim_{\mathrm{fr}}(A_{a,\mathbf S})-\cdim^{\mathrm{impl}}(A_{a,\mathbf S})
=
a+r-\frac12.
\]
特别地，上述裂隙可任意大。
\end{corollary}
\begin{proof}
群 \(A_{a,\mathbf S}\) 是可数阿贝尔群；当 \(a+r>0\) 时它非平凡且无限，因此存在集合单射
\[
A_{a,\mathbf S}\hookrightarrow \NN,
\]
故
\[
\cdim^{\mathrm{impl}}(A_{a,\mathbf S})\le \frac12.
\]
另一方面，\(\cdim^{\mathrm{impl}}(A_{a,\mathbf S})=0\) 意味着 \(A_{a,\mathbf S}\) 可单射入 \(\NN^0\)，即单点集，矛盾。故
\[
\cdim^{\mathrm{impl}}(A_{a,\mathbf S})=\frac12.
\]

再由定理 \ref{thm:cdim-finite-localization-directsum-prime-ledger-exact-sequence} 以及 \(\cdim_{\mathrm{fr}}(\ZZ^{a})=a\)，得到
\[
\cdim_{\mathrm{fr}}(A_{a,\mathbf S})
=
\cdim_{\mathrm{fr}}(\ZZ^a)+\cdim_{\mathrm{fr}}(G_{\mathbf S})
=
a+r.
\]
相减即得所示精确裂隙公式。证毕。
\end{proof}

\leanverified{paper\_cdim\_finite\_localization\_solenoid\_quotient\_embedding\_rigidity}
\begin{theorem}[有限局部化 solenoid 的商端可达与嵌入刚性]\label{thm:cdim-finite-localization-solenoid-quotient-embedding-rigidity}
设 \(\mathbf S=(S_1,\dots,S_r)\) 为有限素数集族，记
\[
G_{\mathbf S}:=\bigoplus_{i=1}^{r}\ZZ[S_i^{-1}],
\qquad
\Sigma_{\mathbf S}:=\widehat{G_{\mathbf S}}.
\]
则成立：
\begin{enumerate}
  \item 存在规范满同态
  \[
  \Sigma_{\mathbf S}\twoheadrightarrow \TT^r,
  \]
  其核同构于
  \[
  \prod_{p\in\mathbb P}\ZZ_p^{\,m_p(\mathbf S)}.
  \]
  \item 若某个 \(S_i\neq\varnothing\)，则对任意有限 \(N\) 都不存在连续单射群同态
  \[
  \Sigma_{\mathbf S}\hookrightarrow \TT^N.
  \]
  \item 因而
  \[
  \Sigma_{\mathbf S}\cong \TT^r
  \iff
  S_1=\cdots=S_r=\varnothing.
  \]
\end{enumerate}
\end{theorem}
\begin{proof}
第 \(1\) 条正是定理 \ref{thm:cdim-finite-localization-directsum-prime-ledger-exact-sequence} 的对偶短正合列。

现证第 \(2\) 条。设某个 \(S_i\neq\varnothing\)，并反设存在连续单射
\[
\iota:\Sigma_{\mathbf S}\hookrightarrow \TT^N.
\]
Pontryagin 对偶把紧交换群的单射送到离散交换群的满射，故得到满射
\[
\iota^\vee:\ZZ^N\twoheadrightarrow G_{\mathbf S}.
\]
然而 \(\ZZ^N\) 为有限生成阿贝尔群，故其任意商群仍有限生成；另一方面，由定理
\ref{thm:cdim-finite-localization-directsum-prime-ledger-exact-sequence}，
\[
G_{\mathbf S}/\ZZ^r
\cong
\bigoplus_{p\in\mathbb P}C_{p^\infty}^{\,m_p(\mathbf S)}.
\]
若某个 \(S_i\neq\varnothing\)，则存在素数 \(p\) 使 \(m_p(\mathbf S)\ge 1\)，从而 \(G_{\mathbf S}\) 具有非零 Prüfer 商因子，特别地不是有限生成群。这与 \(\iota^\vee\) 的满射性矛盾。故所求单射不存在。

最后，若所有 \(S_i=\varnothing\)，则
\[
G_{\mathbf S}\cong \ZZ^r,
\qquad
\Sigma_{\mathbf S}\cong \TT^r.
\]
反之，若某个 \(S_i\neq\varnothing\)，则第 \(2\) 条表明 \(\Sigma_{\mathbf S}\) 不可能同构于任何有限维环面，尤其不可能同构于 \(\TT^r\)。证毕。
\end{proof}

\noindent
由此，有限局部化对象族在本节的圆维语境中获得了完全闭合的双轴图景：离散侧由有限秩自由部分与 Prüfer 素数账本的精确短正合列刻画，紧侧则表现为环面商与 profinite 核的规范扩张；这一结构既排除了把全部自然数乘法压入任意有限秩局部化加性账本的可能性，也使实现层最小编码与结构层有限秩圆维之间的裂隙可以被精确计算，并在 solenoid 一侧外显为“可商向有限维环面而不可反向嵌入”的严格不对称性。

\endinput


\begin{theorem}[素幂粗分辨率的有限商增长谱与素数多重度反演]\label{thm:cdim-star-mixed-localization-pk-growth-inversion}
\leanverified{paper\_cdim\_star\_mixed\_localization\_pk\_growth\_inversion}
在定理 \ref{thm:cdim-star-mixed-localization-two-axis-decomposition} 的设定下，固定素数 \(p\) 并对 \(k\ge 1\) 定义
\[
Q_p(k):=\card{A/p^{k}A}.
\]
则对一切 \(k\ge 1\) 有
\[
Q_p(k)=p^{k(r-m_p)}.
\]
因此 \(m_p\) 可由有限商增长精确恢复：
\[
m_p=r-\frac{1}{k}\log_{p}Q_p(k)\qquad(\forall\,k\ge 1).
\]
并且
\[
\max_{p\in\mathbb P}m_p
=
r-\min_{p\in\mathbb P}\Bigl(\frac{1}{k}\log_p\card{A/p^kA}\Bigr)\qquad(\forall\,k\ge 1).
\]
\end{theorem}
\begin{proof}
由 \(A=\bigoplus_{i=1}^{r}A_i\) 得 \(A/p^{k}A\cong\bigoplus_{i=1}^{r}A_i/p^{k}A_i\)。
对单个分量 \(A_i=\ZZ[1/S_i]\) 分两类讨论。
\par
若 \(p\in S_i\)，则 \(p\) 在 \(A_i\) 中可逆，乘法映射 \([p^k]_{A_i}\) 为双射，故 \(A_i/p^{k}A_i=0\)。
若 \(p\notin S_i\)，则对任意 \(t\ge 0\) 有 \(\gcd(p^k,\,(\prod_{q\in S_i}q)^t)=1\)，从而存在自然满射
\[
\phi_i:A_i\to \ZZ/p^{k}\ZZ,
\qquad
\phi_i\!\left(\frac{a}{b}\right):=a\cdot(b^{-1}\bmod p^{k}),
\]
其中 \(b\) 的素因子全在 \(S_i\)。
其核恰为 \(p^{k}A_i\)，故 \(A_i/p^{k}A_i\cong \ZZ/p^{k}\ZZ\)。
\par
于是
\[
\card{A/p^{k}A}
=\prod_{i:\,p\notin S_i}\card{A_i/p^{k}A_i}
=\prod_{i:\,p\notin S_i}p^{k}
=p^{k\cdot \card{\{\,i:\ p\notin S_i\,\}}}
=p^{k(r-m_p)}.
\]
其余等式由代数化简直接得到。证毕。
\end{proof}

\begin{proposition}[有限素数截断映射的核]\label{prop:cdim-star-finite-prime-truncation-kernels}
\leanverified{paper\_cdim\_star\_finite\_prime\_truncation\_kernels}
若 \(S\subseteq S'\) 为有限素数集，则包含 \(A_S\hookrightarrow A_{S'}\) 诱导自然满同态
\[
\pi_{S',S}:\widehat{A_{S'}}\twoheadrightarrow \widehat{A_S},
\qquad
\ker(\pi_{S',S})\cong \prod_{p\in S'\setminus S}\ZZ_p.
\]
并且对任意有限 \(S\) 有自然满同态
\[
\pi_S:\widehat{\QQ}\twoheadrightarrow \widehat{A_S},
\qquad
\ker(\pi_S)\cong \prod_{p\notin S}\ZZ_p.
\]
\end{proposition}
\begin{proof}
对偶化等价于计算离散侧商群：
\[
A_{S'}/A_S\cong \bigoplus_{p\in S'\setminus S}\ZZ(p^\infty),\qquad
\QQ/A_S\cong \bigoplus_{p\notin S}\ZZ(p^\infty).
\]
再对偶化并用
\(
\widehat{\bigoplus_i B_i}\cong \prod_i \widehat{B_i}
\)
即得结论。证毕。
\end{proof}

\begin{corollary}[有限素数截断不存在连续分裂截面]\label{cor:cdim-star-no-continuous-section-finite-truncation}
\leanverified{paper\_cdim\_star\_no\_continuous\_section\_finite\_truncation}
令 \(S\subsetneq\mathbb P\) 为有限素数集，并令
\(
\pi_S:\widehat{\QQ}\twoheadrightarrow K_S=\widehat{\ZZ[1/S]}
\)
为命题 \ref{prop:cdim-star-finite-prime-truncation-kernels} 中的标准满同态。
则不存在连续群同态 \(g:K_S\to \widehat{\QQ}\) 使
\[
\pi_S\circ g=\mathrm{id}_{K_S}.
\]
等价地，短正合列
\[
0\to \prod_{p\notin S}\ZZ_p\to \widehat{\QQ}\xrightarrow{\ \pi_S\ }K_S\to 0
\]
在拓扑群范畴中不分裂。
\end{corollary}

\begin{proof}
反设存在连续群同态 \(g\) 使 \(\pi_S\circ g=\mathrm{id}\)。
对偶化得到离散侧的群同态
\[
\widehat g:\QQ\to \ZZ[1/S]
\]
满足 \(\widehat g|_{\ZZ[1/S]}=\mathrm{id}_{\ZZ[1/S]}\)。
取任意素数 \(p\notin S\)。
则 \(1/p\in\QQ\)，并且由于 \(\widehat g\) 为群同态且在 \(\ZZ\) 上恒等，有
\[
p\cdot \widehat g(1/p)=\widehat g(1)=1,
\]
从而 \(\widehat g(1/p)=1/p\in \ZZ[1/S]\)，与 \(p\notin S\) 矛盾。
故 \(g\) 不存在。证毕。
\end{proof}

\begin{theorem}[局部化 solenoid 在 \texorpdfstring{$\TT$}{T} 上的态射分类]\label{thm:cdim-solenoid-over-t-morphism-classification}
\leanverified{paper\_cdim\_solenoid\_over\_t\_morphism\_classification}
令 \(S_1,S_2\subset\mathbb{P}\) 为有限素数集，并记
\[
K_{S_i}:=\widehat{\ZZ[1/S_i]},\qquad \pi_{S_i}:K_{S_i}\twoheadrightarrow\TT
\]
为对偶于 \(\ZZ\hookrightarrow \ZZ[1/S_i]\) 的标准投影。
考虑满足 \(\pi_{S_2}\circ f=\pi_{S_1}\) 的连续群同态 \(f:K_{S_1}\to K_{S_2}\)。
则：
\begin{enumerate}
  \item 此类 \(f\) 存在当且仅当 \(S_2\subseteq S_1\)；
  \item 当 \(S_2\subseteq S_1\) 时，此类 \(f\) 唯一且为满射，并且
  \[
  \ker(f)\cong \prod_{p\in S_1\setminus S_2}\ZZ_p.
  \]
  特别地，对任意素数 \(p\)，\(\ker(f)\) 的 Sylow pro-\(p\) 分量非平凡当且仅当 \(p\in S_1\setminus S_2\)，从而差集 \(S_1\setminus S_2\) 可由 \(\ker(f)\) 的素分量结构完全恢复。
\end{enumerate}
\end{theorem}
\begin{proof}
把条件 \(\pi_{S_2}\circ f=\pi_{S_1}\) 对偶化。
由 Pontryagin 对偶的反变等价，\(f\) 等价于离散侧的群同态
\[
\widehat f:\ZZ[1/S_2]\longrightarrow \ZZ[1/S_1]
\]
满足 \(\widehat f|_{\ZZ}=\mathrm{id}_{\ZZ}\)。
若 \(p\in S_2\)，则 \(1/p\in \ZZ[1/S_2]\)，且由 \(\ZZ\)-线性有
\(
p\cdot \widehat f(1/p)=\widehat f(1)=1
\)，
故 \(\widehat f(1/p)=1/p\in \ZZ[1/S_1]\)，从而 \(p\in S_1\)。
因此存在性推出 \(S_2\subseteq S_1\)。
\par
若 \(S_2\subseteq S_1\)，则自然包含 \(\ZZ[1/S_2]\hookrightarrow \ZZ[1/S_1]\) 满足 \(\widehat f|_{\ZZ}=\mathrm{id}\)，从而对偶得到所需 \(f\)；由对偶化的唯一性，该 \(f\) 唯一。
并且 \(\widehat f\) 为单射，故 \(f\) 为满射。
\par
核结构由命题 \ref{prop:cdim-star-finite-prime-truncation-kernels} 取 \(S=S_2\)、\(S'=S_1\) 直接给出：
\(
\ker(f)\cong\prod_{p\in S_1\setminus S_2}\ZZ_p
\)。
Sylow 分量的刻画由 \(\ZZ_p\) 的 pro-\(p\) 性与素数坐标正交性（定理 \ref{thm:cdim-star-prime-coordinate-orthogonality}）立即推出。证毕。
\end{proof}

\begin{theorem}[素数坐标的连续正交性]\label{thm:cdim-star-prime-coordinate-orthogonality}
\leanverified{paper\_cdim\_star\_prime\_coordinate\_orthogonality}
对任意互异素数 \(p\neq q\)，有
\[
\operatorname{Hom}_{\mathrm{cont}}(\ZZ_p,\ZZ_q)=0.
\]
因此，对任意有限集 \(S,T\)，每个连续同态
\[
f:\prod_{p\in S}\ZZ_p\longrightarrow \prod_{q\in T}\ZZ_q
\]
都按坐标分解，且仅在相同素数坐标上可能非零。
\end{theorem}
\begin{proof}
\(\ZZ_p\) 为 pro-\(p\) 群，其连续像仍为 pro-\(p\) 群；
\(\ZZ_q\) 的任意有限连续商群阶为 \(q^n\)，故不含非平凡 pro-\(p\) 闭子群（\(p\neq q\)）。
于是任意连续同态 \(\ZZ_p\to\ZZ_q\) 只能为零映射。

对乘积同态 \(f\)，令
\[
f_{q,p}:=\operatorname{pr}_q\circ f\circ \operatorname{in}_p:\ZZ_p\to\ZZ_q.
\]
若 \(p\neq q\)，则 \(f_{q,p}=0\)。
因此 \(f\) 由对角块 \(f_{p,p}\) 唯一决定。证毕。
\end{proof}

\begin{proposition}[通用 solenoid 的相位--素数纤维模型]\label{prop:cdim-star-universal-solenoid-fiber-model}
在 \(\widehat{\ZZ}\times \RR\) 上定义 \(\ZZ\)-对角作用
\[
n\cdot(x,t)=(x+n,t-n),\qquad n\in\ZZ.
\]
则有同构
\[
\widehat{\QQ}\cong (\widehat{\ZZ}\times \RR)/\ZZ
\cong \widehat{\ZZ}\times_{\ZZ}\RR,
\]
并与 \eqref{eq:pom-universal-solenoid-quotient}--\eqref{eq:pom-universal-solenoid-ses} 的模型一致。
\end{proposition}
\begin{proof}
由备注 \ref{rem:pom-solenoid-extension} 的定义，\(\Sol=(\RR\times\widehat{\ZZ})/\ZZ\) 且 \(\Sol\cong\widehat{\QQ}\)。
交换因子顺序即得上述表示。证毕。
\end{proof}
\leanverified{paper\_cdim\_star\_universal\_solenoid\_fiber\_model}

\begin{theorem}[局部化与有理化的自同态度值域]\label{thm:cdim-star-endomorphism-localized-degree}
记
\[
K_S:=\widehat{\ZZ[1/S]},\qquad K:=\widehat{\QQ}.
\]
则存在自然反同构
\[
\operatorname{End}_{\mathrm{cont}}(K_S)\cong \operatorname{End}(\ZZ[1/S])\cong \ZZ[1/S],
\]
\[
\operatorname{End}_{\mathrm{cont}}(K)\cong \operatorname{End}(\QQ)\cong \QQ.
\]
由于 \(\ZZ[1/S]\) 与 \(\QQ\) 交换，上述反同构可直接视作环同构。
\end{theorem}
\begin{proof}
Pontryagin 对偶给出
\(
\operatorname{Hom}_{\mathrm{cont}}(\widehat B,\widehat A)\cong \operatorname{Hom}(A,B)
\)
的反变等价。
因此
\[
\operatorname{End}_{\mathrm{cont}}(K_S)\cong \operatorname{End}(\ZZ[1/S]),\qquad
\operatorname{End}_{\mathrm{cont}}(K)\cong \operatorname{End}(\QQ).
\]
对 \(A_S=\ZZ[1/S]\)，端同态 \(f\) 由 \(f(1)\) 唯一决定，且
\(
f(1/p^n)=f(1)/p^n
\)
迫使 \(f\) 为“乘以某个 \(a\in \ZZ[1/S]\)”；
故 \(\operatorname{End}(\ZZ[1/S])\cong \ZZ[1/S]\)。
\(\QQ\) 情形同理，因其是一维 \(\QQ\)-向量空间，端同态即“乘以 \(q\in\QQ\)”。
证毕。
\end{proof}
\leanverified{paper\_cdim\_star\_endomorphism\_localized\_degree}

\begin{remark}[winding-register 的局部化/有理化]\label{rem:cdim-star-winding-localization}
在 \(\TT\) 上，自同态的度取值于 \(\ZZ\)。
定理 \ref{thm:cdim-star-endomorphism-localized-degree} 表明：当系统提升到 \(K_S\) 或 \(K\)（即进入 prime-register 的 solenoid 几何）后，自同态由对偶侧标量唯一参数化，其值域分别提升为 \(\ZZ[1/S]\) 与 \(\QQ\)。
在本文寄存器语义下，这一标量可作为缠绕度/覆盖度的代数读出。
\end{remark}

\begin{definition}[扩展半圆维]\label{def:cdim-plus-star}
设 \(M\) 为可消交换幺半群，\(G(M)\) 为其 Grothendieck 群完成。
定义
\[
\cdim_+^\ast(M):=\frac12\,\cdim^\ast\!\bigl(G(M)\bigr).
\]
\end{definition}

\begin{proposition}[局部化正锥的半圆维与素数签名]\label{prop:cdim-plus-star-localized-cone}
设
\[
M_S:=\ZZ[1/S]_{\ge 0}.
\]
则
\[
G(M_S)\cong \ZZ[1/S],\qquad
\cdim_+^\ast(M_S)=\frac12.
\]
同时，\(S\) 的差异不体现在数值 \(\cdim_+^\ast\) 上，而体现在纤维签名
\[
\prod_{p\in S}\ZZ_p
\]
上。
\end{proposition}
\begin{proof}
任意 \(x\in \ZZ[1/S]\) 可写成 \(x=x_+-x_-\) 且 \(x_\pm\in M_S\)，故 \(G(M_S)\cong \ZZ[1/S]\)。
由定理 \ref{thm:cdim-star-z1s-dual-extension} 得 \(\cdim^\ast(\ZZ[1/S])=1\)，代入定义 \ref{def:cdim-plus-star} 即得 \(\cdim_+^\ast(M_S)=1/2\)。
纤维签名由同一定理中的短正合列核给出。证毕。
\end{proof}
\leanverified{paper\_cdim\_plus\_star\_localized\_cone}

\begin{corollary}[有限 prime-register 的精确覆盖边界]\label{cor:cdim-star-finite-prime-register-limit}
设 \(A\subseteq \QQ\) 为离散子群。
若存在有限素数集 \(S\) 使 \(A\) 中每个元素分母的素因子都属于 \(S\)，则
\[
A\subseteq \ZZ[1/S].
\]
反之，若 \(A\) 的分母素因子支持为无限集，则不存在有限 \(S\) 使 \(A\subseteq \ZZ[1/S]\)。
因此，任何仅依赖有限 prime-register 纤维
\(
\prod_{p\in S}\ZZ_p
\)
的精确表征，最多覆盖 \(\ZZ[1/S]\) 型对象，不能精确覆盖全体 \(\QQ\)。
\end{corollary}
\begin{proof}
前一断言由 \(\ZZ[1/S]\) 的定义直接得到；
后一断言由反证法：若 \(A\subseteq \ZZ[1/S]\) 且 \(S\) 有限，则 \(A\) 的分母素支持必有限，与假设矛盾。
最后结论由命题 \ref{prop:cdim-star-finite-prime-truncation-kernels} 的核结构与上述包含判据合并得到。证毕。
\end{proof}
\leanverified{paper\_cdim\_star\_finite\_prime\_register\_limit}

\begin{remark}[主分量与签名分量的分离]\label{rem:cdim-star-main-vs-signature}
对任意有限 \(S\)，\(\cdim^\ast(\ZZ[1/S])=1\) 保持不变；
但纤维 \(\prod_{p\in S}\ZZ_p\) 随 \(S\) 改变。
因此“圆维主分量”与“prime-register 签名分量”构成互补不变量：前者刻画连续相位通道数，后者刻画素数纤维结构。
\end{remark}

\endinput
